% Créditos de las fotografías y de las ilustraciones generadas por IA del
% libro 3 (Química universitaria, segundo año), edición española. Los
% registros de licencia legibles por máquina están en
% images/book3/CREDITS.md; mantener ambos archivos sincronizados con la
% versión inglesa (frontmatter/image-credits-book3.tex).
\chapter*{Créditos de las imágenes}

Los dibujos de este libro son originales. Las fotografías, cuando se
usan, se emplean bajo sus respectivas licencias libres, con nuestro
agradecimiento a sus autores; los enlaces completos a las fuentes y las
URL de las licencias de cada fotografía figuran en
\texttt{images/book3/CREDITS.md}, en el repositorio del libro.
Los pictogramas de peligro son los del Sistema Globalmente Armonizado de
clasificación y etiquetado de productos químicos, tal como los dibuja el
paquete \texttt{ghsystem} de \TeX\ Live.

\bigskip
Junto a las fotografías y a los dibujos originales, algunas figuras son
ilustraciones generadas por IA, producidas con la generación de imágenes
de OpenAI a partir de instrucciones escritas por los editores del libro,
y revisadas en cuanto a su exactitud química antes de incluirlas. La
instrucción completa de cada imagen generada está registrada en
\texttt{images/book3/ai/PROMPTS.md}, en el repositorio.

\bigskip
\noindent Fotografías, por capítulo:
\begin{itemize}\small
  \item Cap.~1: Germain Henri Hess, retrato de I.~I. Reimers (1837--1838),
        CC BY-SA 4.0; Marcellin Berthelot, autor desconocido, dominio público
        (ambas de Wikimedia Commons).
  \item Cap.~2: Josiah Willard Gibbs, fotógrafo desconocido, dominio público
        (Wikimedia Commons).
  \item Cap.~3: François-Marie Raoult, fotógrafo desconocido, dominio público
        (Wikimedia Commons).
  \item Cap.~4: Jacobus Henricus van 't Hoff, de Nicola Perscheid (1904),
        dominio público; Henry Le Chatelier, autor desconocido, CC BY-SA 4.0;
        Fritz Haber y Carl Bosch, Fundación Nobel, dominio público (todas de
        Wikimedia Commons).
  \item Cap.~6: alto horno de Duisburg-Nord, de Dietmar Rabich, CC BY-SA
        4.0; soldadura aluminotérmica de un carril, de PetrS., CC BY-SA 3.0
        (ambas de Wikimedia Commons).
  \item Cap.~7: columna de destilación de una refinería, de CEphoto, Uwe Aranas,
        CC BY-SA 3.0 (Wikimedia Commons).
  \item Cap.~9: Michael Faraday, de Thomas Phillips (1842), dominio público
        (Wikimedia Commons).
  \item Cap.~10: Jaroslav Heyrovsk\'y, archivo del Instituto de Química Física
        J.~Heyrovsk\'y, CC BY-SA 3.0 (Wikimedia Commons).
  \item Cap.~11: una pila de Volta, de GuidoB, CC BY-SA 3.0; Charles Martin Hall
        y Paul H\'eroult, autores desconocidos, dominio público (todas de Wikimedia
        Commons).
  \item Cap.~12: una estatua de cobre, de Elcobbola, dominio público; un ánodo de
        sacrificio, de Zwergelstern, CC BY-SA 3.0 (ambas de Wikimedia Commons).
  \item Cap.~13: Erwin Schr\"odinger, Fundación Nobel (1933), dominio público
        (Wikimedia Commons).
  \item Cap.~14: oxígeno líquido en un imán, de Pieter Kuiper, dominio público;
        Robert Mulliken y Friedrich Hund, de GFHund, CC BY 3.0 (ambas de Wikimedia
        Commons).
  \item Cap.~16: August Kekul\'e (1873), autor desconocido, dominio público
        (Wikimedia Commons).
  \item Cap.~18: disoluciones de metales de transición, de Benjah-bmm27, dominio público; Alfred
        Werner, Les Prix Nobel (1914), dominio público (ambas de Wikimedia Commons).
  \item Cap.~19: un cristal de rubí, dominio público; una esmeralda, de Didier Descouens,
        CC BY-SA 3.0; cristales de sulfato de cobre, de Stephanb, CC BY-SA 3.0 (todas de
        Wikimedia Commons).
  \item Cap.~20: Akira Suzuki, Ei-ichi Negishi y Richard Heck, de BloodIce, CC BY-SA
        4.0 (Wikimedia Commons).
  \item Cap.~22: Charles Friedel (década de 1890) y James Mason Crafts (\textit{Popular
        Science Monthly}, 1900), dominio público (Wikimedia Commons).
  \item Cap.~24: Michel Eug\`ene Chevreul, de Nicolas Eustache Maurin, dominio público
        (Wikimedia Commons).
  \item Cap.~25: Alexander Borodin, autor desconocido, dominio público; Ludwig Claisen, de
        Veronica.villagrasa, CC BY-SA 4.0 (ambas de Wikimedia Commons).
  \item Cap.~26: Robert Robinson, autor desconocido, dominio público (Wikimedia Commons).
  \item Cap.~28: Elias James Corey, de Trvthchem, CC0 (Wikimedia Commons).
  \item Cap.~29: Wallace Carothers en el laboratorio, fotógrafo desconocido, dominio público
        (Wikimedia Commons).
  \item Cap.~30: Emil Fischer, Atelier Victoria, Berlín, hacia 1895, dominio público (Wikimedia
        Commons).
  \item Cap.~31: Mijaíl Tsvet, autor desconocido, dominio público (Wikimedia Commons).
  \item Cap.~32: Francis William Aston (1922), autor desconocido, dominio público; Gustav Kirchhoff,
        Robert Bunsen y Henry Roscoe (1862), dominio público (ambas de Wikimedia Commons).
  \item Cap.~34: William Sealy Gosset (1908), dominio público (Wikimedia Commons).
%% next chapter here
\end{itemize}
\clearpage
